
\section*{Agradecimientos}

A veces, una tesis es mucho más que un trabajo académico; es el cierre de una etapa intensa, la prueba de que, a pesar de los desafíos, siempre podemos seguir adelante. Por eso, este agradecimiento no es una simple formalidad, sino un \textbf{intento sincero de honrar a quienes hicieron posible este recorrido}.

\bigskip

A mis \textbf{padres}, gracias por sostenerme en todos los sentidos imaginables. Por crear las condiciones para que pudiera estudiar, crecer, equivocarme y, aun así, perseverar. Su apoyo incondicional ha sido el suelo firme sobre el que pude construir este camino.


\bigskip

A \textbf{Mariana}, mi psicóloga, por acompañarme con una sensibilidad tan particular y valiente. Gracias por ayudarme a reconstruirme cuando fue necesario y por sostener con tanto cuidado mi proceso hacia una versión más completa e íntegra de mí misma.

\bigskip

A mi director, \textbf{Leonardo Makinistian}, por su claridad, su profundo respeto y su disposición constante. Encontrar una guía así en medio de una carrera tan exigente es una verdadera fortuna, y es algo que valoro inmensamente.

\bigskip

A la \textbf{Universidad Nacional de San Luis}, por ser el espacio donde pude estudiar Física sin barreras económicas y por recordarme, día a día, que la educación pública tiene el poder de transformar realidades.

\bigskip

Al \textbf{Departamento de Física}, por ser un lugar donde siempre encontré a alguien dispuesto a ayudar. Esos gestos, a veces tan sutiles, marcan una enorme diferencia en el trayecto.

\bigskip
\bigskip

Y también, me agradezco \textbf{a mí misma}.

\medskip

Por haber llegado hasta aquí, con todo lo que eso conlleva. A este cuerpo y esta mente que siguieron adelante incluso en los días más inciertos. A las decisiones pequeñas, los desvíos invisibles y los esfuerzos silenciosos que, sin saberlo, hicieron posible este momento.

\bigskip

Gracias, por no soltar.





{\em }

%%%% email addresses
\blfootnote{ }

%\pagebreak

%\setcounter{page}0
%\thispagestyle{empty}
\parskip = 0.2in

\pagebreak

